\section{Evaluation Methodology}
\label{sec:evaluation}

We organize the evaluation around four research questions that follow the
paper's evidence chain.  RQ1 asks whether tool calls actually commit compound
or pre-state-dependent effects.  RQ2 asks whether coarse representations lose
policy-relevant distinctions and whether counterfactual registration recovers
them.  RQ3 asks whether concrete atoms can implement a frozen authorization
relation.  RQ4 asks what security, utility, and audit properties a
descriptor-driven runtime monitor provides.  The questions use separate
oracles and are not interchangeable: runtime attack success, for example,
cannot establish representation sufficiency.

\paragraph{RQ1: Do calls commit independently authorizable effects?}
We replay all 339 calls in AgentDojo v1.1.2's 97 official benign trajectories
from fresh prescribed states.  Before/after differences are normalized into
resource, target, visibility, message, membership, and external-interaction
effects.  We count both the number of logical effect occurrences and whether a
call crosses effect types or benchmark subsystems.

For a controlled representation test, we use 56 calls to five source-verified
AgentDojo tools.  Under the power set of observed effect occurrences, unequal
concrete effect multisets are authorization-separable.  Grouping calls by tool
name, common fields, raw arguments, or typed effects therefore exposes exactly
which policy distinctions each representation erases.

\paragraph{RQ2: Can counterfactual registration recover required distinctions?}
For each registered AgentDojo field, the harness attempts observed-value,
typed, boundary, list-expansion, omission/default, and surface-only variants.
It records committed-state change, output-only change, invariance, invalidity,
and unresolved outcomes separately.  A state-changing pair establishes effect
relevance over the exercised intervention family.  It establishes
authorization necessity only when the policy oracle also separates the pair.
The initial intervention suite is used to refine the candidate representation,
not to estimate held-out generalization.

We therefore freeze a second protocol over five previously unused ToolSandbox
tools and 32 call/state contexts~\cite{lu-etal-2025-toolsandbox}.  These contexts
exercise values, preconditions, defaults, and compound effects; invalid
contexts remain in the results.  We compare the source-effect partition with
partitions induced by common fields and frozen typed atoms.  The provenance
component used later at runtime comes from trusted benchmark evidence
boundaries and is evaluated as a policy input rather than source-validated as a
tool semantic.

\paragraph{RQ3: Can atoms realize the frozen policy relation?}
The 32 ToolSandbox contexts also define a finite mechanism test.  Within each
tool, every context serves as an authority anchor and is paired with every
candidate context, producing 232 ordered queries.  The source oracle allows a
candidate exactly when its executed effect multiset is a submultiset of the
anchor's.  We compare tool identity, canonical exact arguments, common effect
fields, concrete typed atoms, and the source oracle.  This isolates whether the
atom interface can implement the specified relation; the anchors are finite
source-derived bounds, not deployment ACLs.

\paragraph{RQ4: What does descriptor-driven runtime mediation establish?}
AgentDojo supplies 97 benign tasks and 629 official
\texttt{important\_instructions} attack pairs across workspace, Slack, travel,
and banking.  Under DeepSeek, we compare no guard, Spotlighting, and \sys{} with
the same model, decoding, benchmark version, task keys, and validators.  Benign
utility is repeated four times per method in interleaved order.  A second
common-protocol comparison runs no guard, Spotlighting, Prompt Sandwiching, a
PromptArmor-style local adapter, and \sys{} on all 726 keys under one local
Qwen3-32B checkpoint.  These adapters are comparable implementations, not
original-paper reproductions.

Generalization tests use a frozen 320-case manifest spanning four suites and
four public attack templates, plus a separate 40-case bounded-search
diagnostic.  Neither permits result-informed repair.  For cross-environment
transfer, we replay 303 fixed public AgentLAB attack artifacts under AgentDojo
v1.2.1~\cite{agentlab}.  Matched no-guard and provenance-normalized \sys{}
conditions use the same Qwen3-32B checkpoint and case manifest.  This transfer
tests saved evidence under a new harness; it does not reproduce AgentLAB's
cumulative attack optimization.

We separate runtime security from mechanism attribution.  A retrospective
diagnostic applies four views to the same completed no-guard trajectories and
measures only whether a view could intercept a call.  A closed-loop 321-case
applicability subset instead reruns no blocking, whole-call provenance,
effect-only checking, \sys{}, and a raw-field taint control.  The control holds
task grounding, provenance, totalization, and execution fixed while removing
descriptor effect-label expansion.  This comparison tests whether registered
effect semantics change decisions; source-executed collisions remain the
independent evidence for representation quality because the AgentDojo registry
retains every schema field.

\paragraph{Metrics and integrity.}
Official attack success and task utility use AgentDojo's native validators; no
LLM judges either outcome.  We also report coverage, abstention, direct denials,
runtime LLM calls, and executions without a preceding allow.  Exact-key
security comparisons use McNemar tests.  Benign utility uses task-clustered
bootstrap intervals, with non-inferiority requiring the one-sided 95\% lower
bound to exceed $-0.05$.  The finite authorizer reports unsafe pre-allow, false
denial, coverage, abstention, and accuracy over all 232 queries.  Every protocol
freezes its model, descriptor, case manifest, and monitor identity; failed or
missing rows remain visible, and no call reaches an external service.
Appendix~\ref{app:open-science} gives the fail-fast reproduction checks.
